\section{Experimental setup}
\label{sec:setup}

\paragraph{Hardware target.} NUCLEO-F439ZI (Cortex-M4 @ \SI{180}{\mega\hertz}, \SI{256}{\kilo\byte} SRAM =
\SI{192}{\kilo\byte} + \SI{64}{\kilo\byte} \gls{CCM}, no NPU). Latencies are measured with the DWT cycle
counter; the static footprint is read from the binary and completed by the stack peak
(\S\ref{sec:method}). The latency budget is \SI{100}{\milli\second} per inference + update cycle (Gap~2).
Samples travel over the \gls{UART} link frame by frame, each frame protected by a CRC: a zero error rate is the
condition for the compared predictions to be the board's own.

\paragraph{Datasets and CL scenarios.} Two industrial datasets with contrasting scenarios:
\emph{Pronostia} (FEMTO bearings~\cite{Nectoux2012Pronostia}, \emph{class-incremental} by normal/fault
condition) and \emph{Monitoring} (industrial equipment, \emph{domain-incremental} by machine type). Two
feature conditions are considered: \texttt{5feat} (embedded subset) and \texttt{all} (native features). On
Monitoring, \texttt{5feat} and \texttt{all} coincide (4 native features).

\paragraph{Energy bench.} Current is recorded with an X-NUCLEO-LPM01A probe, the board being powered by
the probe itself. Campaigns are run at an imposed inference rate over a fixed window, with repetitions, so
that the gap between two cells is attributable to the model and not to the host. A clock-frequency sweep
completes the setup, the peripheral bus frequency being held constant so as not to shift the serial link
rate.

\paragraph{Provenance of the measurements.} Table~\ref{tab:sources} lists the campaigns and, for each, its
nature. This distinction governs the reading of the whole article: a value \emph{measured on-board} and a
value \emph{emulated on PC} do not differ in reliability --- the emulator reproduces the C kernel's
arithmetic exactly --- but they do not answer the same question.

\begin{table}[t]
\centering
\caption{Provenance and nature of the reported measurements.}
\label{tab:sources}
\begin{tabular}{llc}
\toprule
Campaign & Object & Nature \\
\midrule
\texttt{exp\_S36}            & FP32 parity, INT8 legacy, latency, RAM & measured board \\
\texttt{exp\_S39\_ablation}  & Ablation ladder, causal diagnosis      & emulated PC \\
\texttt{exp\_S40\_board\_v2} & Recovery, calibrated v2 kernel         & measured board \\
\texttt{exp\_S46\_board}     & Quantization moment, A/B               & measured board \\
\texttt{exp\_S47\_depth}     & Depth and granularity                  & emulated PC \\
\texttt{exp\_S48\_summary}   & Bit-packing, actual \texttt{.bss}      & measured board \\
\texttt{exp\_S49\_ram}       & Total RAM by component                 & measured board \\
\texttt{exp\_S50\_int8\_latency} & Latency by kernel segment          & measured board \\
\texttt{exp\_S53\_*}         & Energy, rate, frequency                & measured board \\
\bottomrule
\end{tabular}
\end{table}

\paragraph{Identical conditions and honesty rule.} All comparisons share seed, order and normalization
(\S\ref{sec:method}). The board-v2 campaign covers the four per-channel cells (two datasets $\times$ two
protocols); the \texttt{q15} format was not streamed on-board. Any unmeasured quantity is reported as "to
be measured", with no invented value and no substituted estimate; the full registry of these cells,
with the reason for each, is given in Appendix~\ref{app:missing}. Acronyms used in the article are
expanded at first use and collected in the list of acronyms, page~\pageref{sec:acronyms}.
